

\documentclass[../../../include/open-logic-chapter]{subfiles}

\begin{document}

\olchapter{lecture1}{lecture1}{Lecture I}

\olsection{Introduction}

\subsection{Meaning of Probability in Higher Dimensions}

\scalebox{1.5}{%
  $ x \sim \mathcal{N}(\sigma, \sum )$
}
All of these values are vectors and matrices in high dimensions.

\scalebox{1.5}{%
  $ x^{(i)} \sim \mathcal{N}(\sigma^{(i)}, \sum^{(i,i)} ) \text{ and  }
 \mathrm{Cov} ( x^{(i)},x^{(i)}) = \sum^{(i,j)}$
}

But we will be dealing only with $\textbf{Isotropic}$ Gaussian matrices.

$\Large \sum = {\sigma}^{2}\mathcal{I} = \begin{bmatrix} {\sigma}^{2} & 0 & 0 & \dots  & 0 \\
    0 & {\sigma}^{2} & 0 & \dots  & 0 \\ \vdots & \vdots & \vdots & \ddots & \vdots \\ 0 & 0 & 0 & \dots  & {\sigma}^{2} \end{bmatrix} $

This is to simplify the math so that the \textbf{variance} is same in all directions.

\subsection{Mathematical formalism of the forward Process}

\scalebox{1.5}{%
  $ q(x^t \mid x^{t-1}) = \mathcal{N}(\sqrt{ (1 - {\beta}_t)} ( x_{t - 1} ), {\beta}_t \mathrm{I})$

}

The following needs more explanation which may appear later in the document.

The following is just parameterization.

\scalebox{1.5}{%
  $ x_t =\underbrace{(\sqrt{ (1 - {\beta}_t)}}_{{\alpha}_t} ( x_{t - 1} ) + \underbrace{\sqrt{ {\beta}_t }}_{{\alpha}_t} {\epsilon_t}$}

Multiplying we get


\scalebox{1.5}{%
$ (\sqrt{{\alpha}_t {\alpha}_{t-1}})  x_{t-2}  + \sqrt{{\alpha}_t (1 -{\alpha}_{t-1})} {\epsilon}_{t-1} + \sqrt{(1 -{\alpha}_t)} {\epsilon}_t$
}


It appears that two random variables drawn from the same Gaussian are combined to get a variance like this.


The two quantities are $\textbf{ squared }$ and added.


\scalebox{1.5}{%
${\alpha}_t (1 -{\alpha}_{t-1}) + (1 -{\alpha}_t) = \underbrace{1 -{\alpha}_t {\alpha}_{t-1}}_\text{Variance}$
}
\vskip 3ex
Since we combine random variables like this

\vskip 3ex
we get
\vskip 3ex
\scalebox{1.5}{%
$(\sqrt{{\alpha}_t {\alpha}_{t-1}})  x_{t-2}  + \overbrace{(\sqrt{1 -{\alpha}_t {\alpha}_{t-1}})}^\text{Stanard Deviation of the Variance shown above} \epsilon $
}
\vskip 3ex
Therefore we can have a noisy image at step $ \textbf{t} $ as a function of
the clean image at step $\textbf{ 0  }$

\begin{tcolorbox}[fontupper=\Large,sharp corners, colback=teal!20, colframe=Aquamarine!10!blue]
$q(x^t \mid x^{t-1}) = \mathcal{N}(\sqrt{{\bar{\alpha}_t}} x_0 , 1 - {\bar{\alpha}_t}}{\mathrm{I}})$

with


${\bar{\alpha}_t}} = \prod_{i=1}^t {\alpha}_i \text{ and }  {\alpha}_t} = 1 - {\beta}_t$

\end{tcolorbox}

\begin{tcolorbox}[fontupper=\Large,sharp corners, colback=teal!20, colframe=Aquamarine!10!blue]
  Variance of sum of independent Gaussians is sum of respective variances
\end{tcolorbox}
\vskip 3ex



\subsection{Objective}
$\Large  \underbrace{\max_{\theta}}_{\substack{\uparrow \\ \text{Find } \theta \text{ that} \\ \text{maximizes}}} \quad \underbrace{\log P_{\theta}}_{\substack{\uparrow \\ \text{log likelihood} \\ \text{under model}}} \quad
\underbrace{(x_0)}_{\substack{\uparrow \\ \text{training} \\ \text{data}}}
$
\end{array}

\subsection{Joint Probability Distribution}

$ \Large P( x_1, x_2 ) = P( x_1 ) x \underbrace{P(x_! \mid x_2)}_\text{Conditional Probability}$

\subsection{Marginalization}

$ P(x_!) = \int P(x_! \mid x_2)$

\vskip 3ex
It appears that we couldn't possibly marginalize over all possible trajectories using $ \log P_{\theta}(x_0) = \log \int P_{\theta}(x_0,x_{1:T}) dx_{1:T}$

\vskip 3ex
This is intractable.

\subsection{Introduction of $\textbf{Evidence Lower Bound}$}

\scalebox{1.5}{%
$= \int \left(\frac{P_{\theta}(x_{0:T}) q(x_{1:T},x_0)}{\underbrace{q(x_{1:T},x_0)}_{\text{Forward process we have defined given the clean image} x_0}}\right)dx_{1:T}$
}

\vskip 3ex

This is the same as the expectation which is written like this.

\scalebox{1.5}{%
$ = \mathrm{E_{\underbrace{ x_{1:T} \sim q(x_{1:T},x_0)}_\text{Sampled}}} \left(\frac{P_{\theta}(x_{0:T})}{q(x_{1:T},x_0)}\right)$
}

\vskip 3ex
\subsubsection{Jensen's Inequality}

\vskip 3ex

\scalebox{1.5}{%
$ = \log \left[\mathrm{E} \left(\frac{P_{\theta}(x_{0:T})}{q(x_{1:T},x_0)}\right)\right] \geq
  \mathrm{E} \left[\log \left(\frac{P_{\theta}(x_{0:T})}{q(x_{1:T},x_0)}\right)\right]$
}

It appears that due to the chain rule the actual equation is like this. This has to be learnt again.

\vskip 3ex
\begin{tcolorbox}[fontupper=\Large,sharp corners, colback=teal!20, colframe=Aquamarine!10!blue]

$ = \mathrm{E}_{x_0 \sim q(x_0)} \left(\log{P_{\theta}(x_0)}\right) \geq   \mathrm{E}_{x_{0:T} \sim q(x_{0:T})}\right)\right] \left[\log \left(\frac{P_{\theta}(x_{0:T})}{q(x_{1:T} \mid x_0)}\right)\right]$
\end{tcolorbox}

\vskip 3ex
This also seems to indicate that we don't sum over all the trajectories but only those sampled from the forward process $\textbf{ q }$.

\subsubsection{KL Divergence}

$\textbf{ELBO} = - \Sigma_{t=2}^T = \mathrm{KL} \left( q(x_{t-1} \mid x_t, x_0) \parallel P_{\theta}(x_{t-1} \mid x_t)\right) + \textbf{ Some other terms }$
\vskip 3ex
\subsubsection{Show Lower Bound is tractable}

$ q(\overbrace{x_{t-1}}^\text{Less Noisy} \mid \underbrace{x_t}_\text{More noisy image}, \underbrace{x_0}_\text{Clean image}) \text{ is tractable }$

The question is this. What is the distribution of the less noisy image given the more noisy image along with the clean image ?

\subsubsection{Bayes' Rule}
\vskip 3ex
\scalebox{1.5}{%
$ \underbrace{\mathrm{P(A \mid B)}}_{\text{Posterior}} = \frac{\mathrm{P(B \mid A) P(B)}}{ \underbrace{P(B)}_{\text{Marginal Likelihood}}}$
}

$q( x_{t-1} \mid x_t, x_0 ) = \frac{q( x_t \mid x_{t-1}, x_0 ) q(x_{t-1} \mid x_0)}{q(x_t \mid x_0)} $

\vskip 3ex

We can write these equations now based on Bayes' rule.

\scalebox{1.5}{%
$q(x_{t-1} \mid x_t, x_0) = \frac{q(x_t \mid x_{t-1}, x_0)  \overbrace{q(x_{t-1} \mid x_0)}^{\text{Tractable}}}{\underbrace{q(x_t \mid x_0)}_{\text{Tractable and fixed (Constant)}}}$
}

\vskip 3ex
$\textbf{q  is markovian}$
\vskip 3ex
$\Large \textbf{ And }  q(x_{t-1} \mid x_t, x_0) \propto q(x_t \mid x_{t-1})q(x_t \mid  x_0)

\vskip 3ex
As far as the right side of the approximation shown above is concerned, the product of the probabilities of two Gaussians is also Gaussian.

\vskip 3ex
So the first term in the $\textbf{ ELBO }$ derivation is tractable.
Now we have to show that the second term is too.

\vskip 3ex

We assume $\Large P_{\theta}$ is Gaussian.
\scalebox{1.5}{%
$P_{\theta}(x_{t-1} \mid x_t) = \mathrm{N}(\mu_{\theta}(x_t,t),\Sigma_{\theta}(x_t,t) )$
}
\vskip 3ex

\begin{tcolorbox}[fontupper=\Large,sharp corners, colback=teal!20, colframe=Aquamarine!10!blue]
$\text{ Justification If } q(x_t \mid x_{t-1}) \text{ is Gaussian, then } q(x_{t-1} \mid x_t) \text{ is approx. Gaussian }$
\end{tcolorbox}
\vskip 3ex
\subsubsection{Final Loss}
\vskip 3ex
\scalebox{1.5}{%
$ \mathcal{L}_{DDPM} (\theta) = \mathrm{E}_{t \sim Unif,x_0 \sim q_0(x_0), \epsilon \sim \mathcal{N}(0,{I}_d)} \left[ \lVert{\epsilon}^{\theta} (\sqrt{{\bar{\alpha}_t}} x_0  + \sqrt{1 - {\bar{\alpha}_t}} \epsilon, t - \epsilon )\rVert_2^2 \right]$
}

This seems to mean that if we have a noisy images and the original noise, we can compute the actual noise that was added.
\end{document}
